
\section{DEFINIÇÃO DO PROBLEMA}
\label{sec:def_problema}

O problema central deste projeto se enuncia da seguinte forma: \emph{"A Riachuelo coloca a peça nas mãos do cliente no melhor momento da jornada e a mantém anônima e travada, obrigando-o a passar pelo caixa para levá-la embora"}. Esse é o instante de maior intenção de compra de toda a jornada, e o comércio eletrônico não o reproduz, pois o cliente já verificou caimento e acabamento sem depender de prazo de entrega ou de troca. A empresa, contudo, exige que ele recorra ao caixa para concluir a aquisição, porque apenas o operador do ponto de pagamento libera o dispositivo antifurto. Dessa exigência decorrem duas perdas. Para o cliente, a compra depende de fila e de atendimento mesmo quando ele já dispõe de meio de pagamento no telefone celular. Para a Riachuelo, a peça circula pelo salão sem que a operação acompanhe sua localização em tempo real ou seu tempo de exposição, e a empresa só volta a registrar o item quando o operador o escaneia para efetivar a venda. A companhia paga a proteção contra o furto com a fluidez da experiência de compra e com a visibilidade sobre o próprio salão de vendas.

O grupo partiu da fila de pagamento como recorte inicial desse problema e buscou uma solução que tornasse a compra mais ágil e autônoma. A proposta se concentrava nas bolachas identificadoras fixadas às peças de vestuário, com a implementação de uma tecnologia que permitisse ao cliente identificar a peça de forma rápida, efetuar o pagamento por meio de um aplicativo e obter a liberação do dispositivo de alarme.

Durante a visita à Riachuelo, o grupo alterou esse desenho. As etiquetas das peças já incorporam tecnologia RFID, que a empresa emprega no controle de inventário. Em vez de desenvolver um mecanismo de identificação próprio, o grupo optou por explorar as funcionalidades e as informações que essas tags oferecem e integrá-las ao produto concebido.

A partir dessa constatação, o grupo ampliou o escopo do projeto para além da redução do tempo de pagamento. O RFID passou a ser tratado também como meio de aprimorar o mapeamento e a disponibilidade das informações referentes às peças presentes na loja. A proposta incorpora um ecossistema de informação que permite ao cliente consultar quais produtos estão disponíveis em determinada unidade e em qual setor encontrá-los, o que facilita a localização das peças, reduz o esforço de busca e aumenta a confiabilidade das informações oferecidas.

O produto proposto integra, portanto, as tags RFID presentes nas peças de vestuário, as bolachas de alarme e um sistema de informação. Essa integração reúne a identificação automática dos produtos, o acesso às informações de disponibilidade e localização e um processo de pagamento mais ágil, e atende tanto à experiência de compra quanto aos processos operacionais da loja.